% !TeX root = ../main.tex
%
\chapter{Lessons Learned and Path Forward}\label{chap:discussion}

\section{Overview}

This chapter synthesizes the findings of the two systems presented in this thesis—\capacity and \incognitos—and reflects on their broader implications for operating system (OS) security in the era of hardware-defined protection.
While these systems address different layers of the trust stack, they share a common goal: reestablishing coherence between the protection and trust boundaries that define modern computing.
Together, they illustrate that security retrofitting is not merely an engineering exercise, but a rethinking of how trust is distributed between hardware and software components.


\section{Lessons Learned}

Several key insights emerge from the exploration of hardware-assisted retrofitting across these two systems.

\hdr{Retrofitting demands principled abstraction design}
Adapting hardware security primitives is not simply a matter of rewriting OS interfaces.
Without coherent abstraction boundaries, new mechanisms risk semantic drift—hardware may enforce guarantees that the OS does not expose or understand.
\capacity highlights the importance of aligning enforcement granularity (capability-based references) with existing abstractions (file descriptors, pointers).


\hdr{Compatibility is both a constraint and an enabler}
While retrofitting inherently limits radical redesign, it ensures deployability in existing ecosystems.
Both \capacity and \incognitos succeed because they extend, rather than replace, established abstractions (\gls{posix}, file descriptors, memory mappings).
This pragmatic stance enables security evolution without ecosystem fragmentation.

\hdr{4. Security must be co-designed with performance and observability.}
Hardware features introduce subtle trade-offs between efficiency, transparency, and trust.
The evaluation of both systems demonstrates that principled use of hardware capabilities can achieve meaningful security benefits with acceptable overheads.

\section{Open Challenges}

Despite these contributions, several open challenges remain in the pursuit of hardware-software co-design for OS security.


\hdr{Evolving hardware trust anchors}
Future processors will likely expose richer isolation primitives—such as fine-grained enclaves or programmable integrity domains.

Integrating such features coherently will require new interfaces that can express multi-level trust relationships.
Integrating such features coherently will require new interfaces that can express multi-level trust relationships.

Both systems implicitly redefine the \gls{tcb}, yet current formal methods lack tools to reason about partial trust or shifting boundaries.
Future work should explore hybrid verification frameworks that combine symbolic reasoning for software with microarchitectural models of hardware-enforced isolation.
Such frameworks could formally prove the correctness of trust realignment strategies, ensuring that retrofitting does not inadvertently introduce new vulnerabilities.

\hdr{Formal reasoning across protection boundaries}
Current verification frameworks often assume a single trusted kernel.
As OS trust becomes distributed across hardware and user space, reasoning about global invariants will demand compositional verification techniques that span layers.

\hdr{Dynamic trust adaptation}
\incognitos hints at a future where trust may shift dynamically at runtime (e.g., obfuscation levels adjusted to threat conditions).
Generalizing this adaptability beyond unikernels could reshape how security policies are enforced in general-purpose systems.

\hdr{Hardware transparency and accountability}
Retrofitting security also depends on the verifiability of hardware behavior.
Emerging security processors must expose trustworthy telemetry and introspection mechanisms to avoid opaque enforcement that undermines accountability.



